\section{Lineære Afbildninger}

Lineære afbildninger tager en vektor fra et vektorrum og afbilder den til en anden vektor i et andet (eller samme) vektorrum.
For at en afbildning \( L : V \rightarrow W \) skal være lineær, skal den opfylde to betingelser for alle vektorer \( v_1, v_2 \in V \) og alle skalarer \( c \in \mathbb{R} \):
\begin{itemize}
    \item Additivitet: \( L(v_1 + v_2) = L(v_1) + L(v_2) \)
    \item Homogenitet: \( L(c \cdot v) = c \cdot L(v) \)
\end{itemize}

En hurtig test kan være at tjekke om \( L(u + c \cdot v) = L(u) + c \cdot L(v) \).

En anden hurtig test er at tjekke om \( L(0) = 0 \). (dog kun som modbevis)

\kilde{Definition 11.0.1}

\subsection{Matricer som lineære afbildninger}

Givet en matrix \( \mathbf{A} \in \mathbb{F}^{m \times n}\), kan vi definere en lineær afbildning \( L_A : \mathbb{R}^n \rightarrow \mathbb{R}^m \) ved:
\[
    L_\mathbf{A}(v) = \mathbf{A} \cdot v
\]
\kilde{Definition 11.1.1}

Eksempel:

\begin{minipage}{0.45\textwidth}
    \[
    \mathbf{A} = \begin{bmatrix}
        1 & 2 \\
        3 & 4
    \end{bmatrix}, \quad v = \begin{bmatrix}
        5 \\
        6
    \end{bmatrix}
    \]
\end{minipage}
\hfill
\begin{minipage}{0.45\textwidth}
    \[
    L_\mathbf{A}(v) = \mathbf{A} \cdot v = \begin{bmatrix}
        1 & 2 \\
        3 & 4
    \end{bmatrix} \cdot \begin{bmatrix}
        5 \\
        6
    \end{bmatrix} = \begin{bmatrix}
        17 \\
        39
    \end{bmatrix}
    \]
\end{minipage}

\bemærkBig{Alle lineære afbildninger mellem endelig-dimensionelle vektorrum kan repræsenteres som matrix-multiplikationer.}

\subsubsection{Billedet af lineær afbildning}
Givet en afbildnings-matrice \( \mathbf{A} \), er billedet af den lineære afbildning \( L_A \) givet ved  \( \colsp(\mathbf{A}) \).
Se sektion \ref{sec:soejlerum}.

\kilde{Lemma 11.1.2}

\subsubsection{Kerne af lineær afbildning}
Givet en afbildnings-matrice \( \mathbf{A} \), er \(\ker(L_A) = \ker(A)\).
Se sektion \ref{sec:kerne}.

\kilde{Definition 11.1.2}

\subsection{Generelle vektorrum}

\subsubsection{Sammensætning af lineære afbildninger}
Hvis \( L_1 : U \rightarrow V \) og \( L_2 : V \rightarrow W \) er lineære afbildninger, så er sammensætningen \( L_2 \circ L_1 : U \rightarrow W \) også en lineær afbildning. 

\kilde{Sætning 11.2.1}

Hvis \( \alpha, \beta, \gamma \) er baser for henholdsvis \( U, V, W \), så gælder:
\[
_\gamma [L_2 \circ L_1]_\alpha = _\gamma [L_2]_\beta \cdot _\beta [L_1]_\alpha
\]

\subsubsection{Kerne af generelle lineære afbildninger}
Kernen af en lineær afbildning \( L : V \rightarrow W \) er mængden af alle vektorer i \( V \), der afbildes til nulvektoren i \( W \):
\[
\ker(L) = \{ v \in V \mid L(v) = 0 \} 
\kildetag{Definition 11.2.1}
\]

\subsubsection{Dimensionssætningen for lineære afbildninger}
For \(L: V_1 \to V_2\)
\[
\dim(\ker(L)) + \dim(\image(L)) = \dim(V_1)
\kildetag{Korollar 11.4.3}
\]

\subsection{Basisskift som lineær afbildning}
Givet vektorrum $V$ og ordnet basis \( \beta = \{ v_1, v_2, \ldots, v_n \} \)
er funktionen \( L_\beta : V \rightarrow \mathbb{R}^n \) defineret ved:
\[
L_\beta(v) \mapsto [v]_\beta
\]
En lineær afbildning.

Funktionen \(\mathbb{F}^n \to V\) defineret ved
\( (c_1, c_2, \ldots, c_n) \mapsto c_1 v_1 + c_2 v_2 + \ldots + c_n v_n\)
er også en lineær afbildning, og er den inverse af \( L_\beta \).

\kilde{Lemma 11.3.1}

\subsection{Find afbildnings-matrix}
For en lineær afbildning \(L: \mathbb{F}^n \to \mathbb{F}^m\) kan afbildnings-matricen \( [L] \) findes ved at anvende \(L\) på standardbasens vektorer \(e_1, e_2, \ldots, e_n\) og sætte resultaterne som søjler i en matrix:
\[
A = \begin{bmatrix}
    L(e_1) & L(e_2) & \ldots & L(e_n)
\end{bmatrix}
\kildetag{Lemma 11.3.2}
\]

\subsubsection{Find afbildningsmatrix of afbilled en vektor}
Givet vektor
\(
    \begin{bmatrix}
        3 \\
        2 \\
        1
    \end{bmatrix}
\)
og lineær afbildning \( L : \mathbb{R}^3 \rightarrow \mathbb{R}^2 \) defineret ved:
\[
    L\begin{bmatrix}
        x_1 \\
        x_2 \\
        x_3
    \end{bmatrix}
    =
    \begin{bmatrix}
        2x_1 + x_2 \\
        -x_2 + 3x_3
    \end{bmatrix}
\]

Først tager vi den ordnede standardbasis for \(\mathbb{R}^3\):
\(
    e_1 = \begin{bmatrix}
        1 \\
        0 \\
        0
    \end{bmatrix},
    e_2 = \begin{bmatrix}
        0 \\
        1 \\
        0
    \end{bmatrix}, 
    e_3 = \begin{bmatrix}
        0 \\
        0 \\
        1
    \end{bmatrix}
\)
Vi afbilleder nu hver af basisvektorene med \(L\):
\[
    L \begin{bmatrix}
        1 \\
        0 \\
        0
    \end{bmatrix} =
    \begin{bmatrix}
        2 \cdot 1 + 0 \\
        -0 + 3 \cdot 0
    \end{bmatrix} =
    \begin{bmatrix}
        2 \\
        0
    \end{bmatrix}
\quad
    L \begin{bmatrix}
        0 \\
        1 \\
        0
    \end{bmatrix} =
    \begin{bmatrix}
        2 \cdot 0 + 1 \\
        -1 + 3 \cdot 0
    \end{bmatrix} =
    \begin{bmatrix}
        1 \\
        -1
    \end{bmatrix}
\quad
    L \begin{bmatrix}
        0 \\
        0 \\
        1
    \end{bmatrix} =
    \begin{bmatrix}
        2 \cdot 0 + 0 \\
        -0 + 3 \cdot 1
    \end{bmatrix} =
    \begin{bmatrix}
        0 \\
        3
    \end{bmatrix}
\]

Vi har nu afbildningsmatricen:
\(
    [L] =
    \begin{bmatrix}
        2 & 1 & 0 \\
        0 & -1 & 3
    \end{bmatrix}
\)

Og vi kan nu finde afbildningen af vores vektor ved at multiplicere afbildningsmatricen med vektoren:
\[
    L
    \begin{bmatrix}
        3 \\
        2 \\
        1
    \end{bmatrix} 
    =
    \begin{bmatrix}
        2 & 1 & 0 \\
        0 & -1 & 3
    \end{bmatrix} \cdot
    \begin{bmatrix}
        3 \\
        2 \\
        1
    \end{bmatrix} =
    \begin{bmatrix}
        2 \cdot 3 + 1 \cdot 2 + 0 \cdot 1 \\
        0 \cdot 3 + (-1) \cdot 2 + 3 \cdot 1
    \end{bmatrix} =
    \begin{bmatrix}
        8 \\
        1
    \end{bmatrix}
\]

\subsection{Afbildningsmatrice med hensyn til basis}
Givet en lineær afbildning \( L : V \rightarrow W \) med ordnede baser \( \beta = \{ v_1, v_2, \ldots, v_n \} \) for \( V \) og \( \gamma = \{ w_1, w_2, \ldots, w_m \} \) for \( W \), kan afbildningsmatricen \( _\gamma [L]_\beta \) findes ved at afbilde hver basisvektor i \( V \) og udtrykke resultatet i basis \( \gamma \):

\[
_\gamma [L]_\beta = \begin{bmatrix}
    [L(v_1)]_\gamma & [L(v_2)]_\gamma & \ldots & [L(v_n)]_\gamma
\end{bmatrix}
\kildetag{Lemma 11.3.3}
\]

\subsubsection{Om afbildningsmatricer med hensyn til baser}

\begin{align*}
    [L(v)]_\gamma = \gamma [L]_\beta \cdot [v]_\beta \\
    _\delta [M \circ L]_\beta = _\delta [M]_\gamma \cdot _\gamma [L]_\beta
\end{align*}
\kilde{Sætning 11.3.4}

\subsubsection{Afbildningsmatrix på samme vektorrum}
Giver lineær afbildning \( L : V \rightarrow V \) med ordnet basis \( \beta = \{ v_1, v_2, \ldots, v_n \} \) og ordnet basis \( \gamma = \{ w_1, w_2, \ldots, w_n \} \):
\[
_\gamma [id_V]_\beta = \begin{bmatrix}
    [v_1]_\gamma & [v_2]_\gamma & \ldots & [v_n]_\gamma
\end{bmatrix}
\kildetag{Lemma 11.3.5}
\]

\paragraph{Om afbildning på samme vektorrum}

Det er bare en anden måde at skrive en basisskift-matrice.

\begin{align*}
    _\delta [id_V]_\gamma \cdot _\gamma [id_V]_\beta = _\delta [id_V]_\beta \\
    _\beta [id_V]_\beta = I_n \\
    (_\gamma [id_V]_\beta)^{-1} = _\beta [id_V]_\gamma
\end{align*}
\kilde{Lemma 11.3.6}

\subsection{Find vektor der afbildes til en given vektor}
Givet lineær afbildning \( L : V \rightarrow W \) og vektor \( w \in W \), kan vi  \( v \in V \), så \( L(v) = w \), ved at løse ligningen:
\[
_\gamma [L]_\beta \cdot [v]_\beta = [w]_\gamma
\]

For vektoren \([v]_\beta\) indsættes blot $(x,y,z, \cdots)$ og løses som et almindeligt ligningssystem.

Når en løsning er fundet vil alle vektorer der afbildes til \( w \) være givet ved:
\[
    v = v_p + v_{ker}
\]
hvor \( v_p \) er den fundne løsning og \( v_{ker} \) er en vilkårlig vektor i kernen af \( L \). 


\subsection{Eksempler på lineære afbildninger}
\begin{itemize}
    \item Sporet af en matrix:
    \begin{align*}
        \mathbb{R}^{n \times n} &\rightarrow \mathbb{R} \\
        A &\mapsto \sum_{i=1}^{n} a_{ii}
    \end{align*}
    \kilde{Eksempel 11.2.2}
    \item En vektor til samme vektor med færre dimensioner:
    \begin{align*}
        L : \mathbb{R}^3 &\rightarrow \mathbb{R}^2 \\
        \begin{bmatrix}
            x_1 \\ x_2 \\ x_3
        \end{bmatrix}
        &\mapsto
        \begin{bmatrix}
            x_1 \\ x_2
        \end{bmatrix}
    \end{align*}
\end{itemize}

\subsection{Find all vektorer der afbiledes til w ved ordnet basis}
Givet lineær afbildning \(L: V_1 \to V_2\) med ordnet basis \(\beta\) for \(V_1\) og \(\gamma\) for \(V_2\), og vektor \(w \in V_2\).

Det er muligt at finde alle vektorer \(v \in V_1\), så \(L(v) = w\), ved at løse ligningen:
\[
_\gamma [L]_\beta \cdot [v]_\beta = [w]_\gamma
\]
For vektoren \([v]_\beta\) indsættes blot \((x_1, x_2, \ldots, x_n)\) og løses som et almindeligt ligningssystem.

Da er det $v$ som er løsningen til $L(v) = w$.

\subsubsection{Kerne ved samme metode}
På samme måde kan kernen for $L$ findes ved at løse ligningen:
\[
_\gamma [L]_\beta \cdot [v]_\beta = 0
\]

Igen er det her $v$ som er løsningen til $L(v) = 0$.